\section{Validation Methodology}
A precision passive circuit minimizes the uncertainty of the test sample, making instrument variability the dominant source of measurement variation \cite{kasper2022electrochemical}. Screen-printed electrodes, in contrast, vary between units and with repeated use \cite{anshori2025ganestat, finvsgar2018reusability}, so an electrode-based comparison would confound electrode variability with instrument performance. A dummy cell was therefore used to evaluate STEI-ZStat against a commercial reference before biosensor validation. It was treated as an instrument-metrology fixture rather than a model of a particular biological medium or analyte. A factory-calibrated PalmSens4 (PALM-PS4.F2.10), capable of EIS measurements from 10~$\mu$Hz to 1~MHz, was used as the commercial reference.

The external dummy cell was designed to fit into the SPE connector of both instruments (Fig.~\ref{fig:system_design}). It implemented a Randles circuit ($R_s + (R_{ct} \parallel C_{dl})$) with $R_s = 560~\Omega$ ($\pm$0.1\%), $R_{ct} = 10$~k$\Omega$ ($\pm$0.02\%), and $C_{dl} = 33$~nF ($\pm$5\%, YAGEO CC0402JRX7R9BB333). The topology was selected because it is a widely used first-order representation of an electrochemical interface. It also matches the commercial verification cell from which the acceptance criteria were adapted \cite{gamry2011universal} and can be realized entirely from discrete components with specified nominal values. Warburg or constant-phase elements were not included because they would add parameters not directly realized by the discrete components and would complicate the ideal-component validation.

The component values were selected to test distinct regions of the EIS response. They yield high- and low-frequency plateaus of approximately 560~$\Omega$ and 10.56~k$\Omega$ with a characteristic frequency of approximately 482~Hz. Thus, the 0.2~Hz--200~kHz sweep captured both resistive plateaus and the main semicircular response. The selected nominal values are also consistent with those of a previously reported study of potentiostat performance \cite{wong2024cyclic}. The nominal component values were used as the analytical reference. However, the capacitor was not independently measured in the assembled circuit and therefore could not serve as a capacitance reference. The cell was received 134~days before measurement, providing a lower bound on the time since reflow. At 0.2~Hz, the computed impedance magnitude of this circuit is within 0.001\% of $R_s + R_{ct}$, so lower frequencies would provide negligible additional information.

\begin{figure}[!b]
\centering
\includegraphics[width=0.95\columnwidth]{fig2_bode_nyquist_palmsens_vs_bstat.png}
\caption{Overlaid and aggregated seven EIS measurements for STEI-ZStat and PalmSens4 in a (a) Bode plot and a (b) Nyquist plot across 55 frequencies. Error bars are omitted because the replicate standard deviation, with a median of 0.024\% in $|Z|$ for STEI-ZStat and 0.026\% for PalmSens4, is smaller than the plot markers; replicate spread is reported in Table~\ref{tab:repeatability}.}
\label{fig:spectra}
\end{figure}

EIS measurements used a nominal excitation of 30 mV$_p$ (21.21 mV$_{rms}$), 0 VDC bias, and 55 logarithmically spaced frequencies swept from 200~kHz to 0.2~Hz. Because of the 12-bit DAC resolution of the AD5941, the actual excitation voltage generated by STEI-ZStat was 30.23 mV$_p$. Because the dummy cell is passive and linear, this 0.8\% excitation difference does not affect impedance. Seven replicate spectra were acquired per instrument.

The instruments used different acquisition strategies. PalmSens4 integrated for the longer of 0.5~s and two excitation periods per frequency, giving a sweep duration of approximately 111~s. STEI-ZStat used frequency-dependent discrete Fourier transform (DFT) integration windows of 0.02--13.67~s and separately measured the load and $R_{CAL}$, giving a sweep duration of approximately 280~s. The longer low-frequency integration time reduced the effective noise bandwidth of STEI-ZStat relative to PalmSens4, and the shorter high-frequency windows increased it. Integration time and noise bandwidth were therefore not matched, and should be considered when comparing precision.

PalmSens4 used automatic current ranging from 10~mA to 1~nA. STEI-ZStat used the 4.7~k$\Omega$ $R_{CAL}$ as the closest available value to the geometric mean of the expected 560--10,560~$\Omega$ range. The installed $R_{CAL}$ measured 4696~$\Omega$ on a handheld meter with a specified accuracy of $\pm$(0.2\% + 3 digits) at 0.001~k$\Omega$ resolution, so the nominal 4.7~k$\Omega$ value was retained in firmware. Because the AD5941 measures impedance relative to $R_{CAL}$, the 1\% tolerance of $R_{CAL}$ limits the absolute accuracy of the measured impedance magnitude; higher-accuracy $R_{CAL}$ characterization was not performed. The combined transimpedance amplifier (TIA) gain was automatically adjusted from 500~$\Omega$ to 15~k$\Omega$ using discrete RTIA and PGA settings.

The validation procedure comprised four stages. Measurement data were first subjected to reactance-sign screening for positive imaginary impedance ($Z'' > 0$), which is inconsistent with the expected capacitive response of the passive Randles cell. Such points were flagged and retained with weight 0.05 rather than removed. Resistive plateaus were evaluated at 200~kHz ($|Z|\approx R_s$) and 0.2~Hz ($|Z|\approx R_s+R_{ct}$). At 200~kHz, residual reactance makes the ideal $|Z|$ approximately 0.10\% above $R_s$; this offset is included in the reported deviation.

Full-band performance was evaluated by complex nonlinear least-squares (CNLS) regression using the Levenberg--Marquardt algorithm in the lmfit library (v1.3.4) in Python \cite{newville2016lmfit}. An ideal capacitor was used rather than a constant-phase element so that the fit tested the ideal-component assumption. Real and imaginary components were fitted simultaneously with residuals weighted by 1/$|Z|$, with flagged points retained at weight 0.05. Two fits were performed.

The first fit fixed $R_s$ and $R_{ct}$ at their nominal values and fitted $C_{dl}$ to characterize spectral deviation from the nominal Randles model. Because $C_{dl}$ was not independently characterized in circuit, the fitted value was treated as a model parameter rather than a traceable capacitance measurement. The fitted impedance spectrum was compared with the measured spectrum, and the median deviation, maximum deviation, and standard deviation (SD) of $|Z|$ and phase were calculated.

The second fit optimized $R_s$, $R_{ct}$, and $C_{dl}$ to assess parameter recovery and cross-instrument equivalence. Parameter deviations were calculated relative to the nominal component values, and replicate variability was summarized using relative standard deviation (\%RSD). The resulting values were then compared with predefined acceptance criteria adapted from the Gamry Universal Dummy Cell operator's manual \cite{gamry2011universal}. The manual specifies absolute acceptance windows of 196--204~$\Omega$, 2.95--3.07~k$\Omega$, and 0.90--1.10~$\mu$F for the fitted parameters of its own verification cell, with nominal values of 200~$\Omega$, 3.01~k$\Omega$, and 1~$\mu$F, respectively. These correspond to $\pm$2\% for the resistive elements and $\pm$10\% for the capacitance. The same percentages were applied to the present dummy cell. These are vendor tolerances rather than a published standard, are not specifications of the PalmSens4 used in this study, and should not be interpreted as standardized accuracy requirements for EIS instrumentation.

Device equivalence was assessed from the second CNLS fit. Because the seven measurements per instrument were acquired in separate blocks rather than paired trials, the samples were treated as independent ($n = 7$ per instrument). Mean differences and 95\% confidence intervals (CI) were calculated using the Welch--Satterthwaite method. Because the population variance is unknown and is estimated from n = 7 replicates, critical values were taken from the Student's t distribution rather than the normal distribution, with Welch--Satterthwaite degrees of freedom, using SciPy (v1.18.0). Equivalence was accepted when the entire 95\% CI remained within the predefined acceptance window. Although a 90\% interval corresponds to two one-sided tests at $\alpha$ = 0.05, the 95\% CI criterion was intentionally conservative. The resulting limits were $\pm$11.20~$\Omega$ for $R_s$, $\pm$200.00~$\Omega$ for $R_{ct}$, and $\pm$3.30~nF for $C_{dl}$, derived from the nominal component values.

Unit-to-unit reproducibility was assessed four weeks after the primary measurements, using three assembled STEI-ZStat boards measured under the same conditions on the same dummy cell, with seven spectra from each board analyzed to match the replicate count of the primary comparison. The microcontroller module was transferred between boards, so this comparison isolates the analog front end, calibration resistor, and board layout.

Because only one unit of each instrument was evaluated in the primary comparison, the equivalence result applies only to the tested STEI-ZStat and PalmSens4 under the specified passive-circuit, excitation, and acquisition conditions. It does not establish unit-to-unit reproducibility or equivalence of the STEI-ZStat design as a whole.
